\documentclass[11pt, a4paper]{report}

\usepackage[utf8]{inputenc}
\usepackage[T1]{fontenc}
\usepackage[french]{babel}

% Appel des packages dans le dossier preambule
\def\packagepath{../../../../../preambule}   % path du package princial

\usepackage{\packagepath/page_layout} % <--- Nouveau
\usepackage{\packagepath/config_info}
\usepackage{\packagepath/cadres_cours}
\usepackage{\packagepath/zones_reponse}
\usepackage{\packagepath/config_ds}
\usepackage{\packagepath/config_maths}

% --- PILOTAGE ---
%\def\visibleornot{visible}    % visible or invisible
\ifdefined\visibleornot\else\def\visibleornot{invisible}\fi


\def\documentpath{./Sources_Latex}
\graphicspath{{./Images/}}


\def\ptsexoA{0}



\begin{document}

\def\level{Premiere}  % Classe à droite
\def\course{Maths}
\def\chaptername{Suites Arithmétiques} % Titre au centre
\def\docname{\quad}        % Identifiant à gauche

\pagestyle{DS_LP}

\emph{Objectif: Modéliser une suite arithmétique et automatiser le calcul de ses termes à l'aide du langage Python}

\subsubsection{Partie 1: Etude Mathématique}

On considère une suite arithmétique $(u_n)$ de premier terme $u_0=-8$ et de raison $r=3$.

\begin{enumerate}
    \item Rappeler la relation de récurrence qui permet de passer d'un terme au suivant pour la suite $(u_n)$.
    \item Calculer les 5 premiers termes de la suite $(u_n)$.
    \item Donner une formule explicite pour le terme général $u_n$ de la suite
    \item Calculer $u_{100}$ à l'aide de la formule explicite.
\end{enumerate}

\subsubsection{Partie 2: Utilisation d'un algorithme}

\begin{easybox}[L'affectation en python]
 
    Pour stocker une valeur, on utilise une \textbf{variable}. L'instruction suivante signifie \og{} mettre la valeur $-8$ dnas la variable nommée $u$ \fg{}.


    \begin{CodePythontexAlt}[TaillePolice = \large, Largeur=17cm,Centre,PremLigne=1]{}
        # Affectation de la valeur -8 à la variable u
        u = -8
        \end{CodePythontexAlt}

        Il est possible de modifier la valeur d'une variable et de reaffecter cette nouvelle valeur à la même variable. L'ancienne valeur est alors écrasée.

        
    \begin{CodePythontexAlt}[TaillePolice = \large, Largeur=17cm,Centre,PremLigne=1]{}
        # La variable u est remplacée par la valeur de u + 3
        u = u + 3
            
    \end{CodePythontexAlt}
\end{easybox}


\begin{easybox}[Les listes en python]
 
\begin{itemize}
    \item La \textbf{liste}: c'est une structure de données qui permet de stocker plusieurs valeurs dans une même variable. Les éléments d'une liste sont ordonnés et peuvent être de types différents.

    \begin{CodePythontexAlt}[TaillePolice = \large, Largeur=16cm,Centre,PremLigne=1]{}
# Création d'une liste vide
ma_liste = []

# Création d'une liste avec des éléments
ma_liste = [2, 5, 8, 11]

# Ajout d'un élément à la fin de la liste
ma_liste.append(14)
# la liste devient [2, 5, 8, 11, 14]


    \end{CodePythontexAlt}


    \medskip

    \item La \textbf{boucle for}: Pour répéter une action on utilise une boucle:
    \begin{CodePythontexAlt}[TaillePolice = \large, Largeur=16cm,Centre,PremLigne=1]{}
# la boucle est effectuée 10 fois (de 0 à 9)
for i in range(10):
    # affichage de la valeur de i à chaque itération
    print(i)    

\end{CodePythontexAlt}

\end{itemize}


\end{easybox}

\pagebreak


\begin{enumerate}
    \item Compléter le code suivant afin qu'il génère les 5 premiers termes de la suite et les stocke dans une liste nommée \verb|liste_termes|.
    

    \begin{CodePythontexAlt}[TaillePolice = \large, Largeur=17cm,Centre,PremLigne=1]{}
# Initialisation de la liste pour stocker les termes de la suite
liste_termes = ............
# Premier terme de la suite
u = ............
# Raison de la suite
r = ............

# Boucle pour calculer les 5 premiers termes de la suite
for i in range(............):
    # Ajouter le terme actuel à la liste des termes
    ............
    # Calculer le terme suivant en utilisant la relation de récurrence
    ............

print(liste_termes)  # Affiche les 5 premiers termes de la suite
    \end{CodePythontexAlt}
    
    \medskip

    \item Tester votre code en le recopiant dans un environnement Python.


\item Vérifier que les termes affichés dans la console correspondent bien à ceux calculés à la partie 1.
\end{enumerate}

\subsubsection{Partie 3: Pour aller plus loin}

\begin{easybox}[Utilisation des valeurs d'une liste]
    Il est possible d'accéder à un élément d'une liste en utilisant son indice. L'indice du premier élément d'une liste est 0, celui du deuxième élément est 1, etc.    
    \begin{CodePythontexAlt}[TaillePolice = \large, Largeur=17cm,Centre,PremLigne=1]{}
# Création d'une liste de termes
liste_termes = [-8, -5, -2, 1, 4]

# Accès à l'élément d'indice 0 (le premier élément de la liste)
print(liste_termes[0])  # Affiche -8

# Accès à l'élément d'indice 2 (le troisième élément de la liste)
print(liste_termes[2])  # Affiche -2
    \end{CodePythontexAlt}
    
\end{easybox}
\begin{enumerate}
    \item En vous inspirant du code précédent, compléter votre code afin de calculer la somme des termes contenus dans la liste que vous venez de créer.

\begin{CodePythontexAlt}[TaillePolice = \large, Largeur=17cm,Centre,PremLigne=1]{}
# Initialisation de la variable pour stocker la somme
somme = ............

# Boucle pour parcourir les termes de la liste et les additionner
for i in range(............):
    # Ajouter le terme actuel à la somme
    somme = ........................

print(somme)  # Affiche la somme des termes de la suite
\end{CodePythontexAlt}
    
\medskip

\item Calculer la somme des 5 premiers termes de la suite à l'aide de la formule de la somme d'une suite arithmétique et vérifier que le résultat est le même que celui obtenu à l'aide du code python.
\end{enumerate}
\end{document}